\documentclass[11pt,a4paper]{article}
\usepackage[utf8]{inputenc}
\usepackage[T1]{fontenc}
\usepackage{amsmath,amssymb,amsthm}
\usepackage{mathtools}
\usepackage{physics}
\usepackage[margin=2.5cm]{geometry}
\usepackage{hyperref}
\usepackage{cleveref}
\usepackage{enumitem}
\usepackage{xcolor}
\usepackage{tcolorbox}
\usepackage{booktabs}
\usepackage{fancyhdr}
\usepackage{graphicx}

% Custom commands (matching NT-1/NT-2/NT-4/INF-1 conventions)
\newcommand{\Id}{\mathbb{1}}
\newcommand{\Ricsq}{R_{\mu\nu}R^{\mu\nu}}
\newcommand{\Rsq}{R^2}
\newcommand{\Weylsq}{C_{\mu\nu\rho\sigma}C^{\mu\nu\rho\sigma}}
\renewcommand{\dd}{\mathrm{d}}
\newcommand{\ie}{\textit{i.e.}}
\newcommand{\eg}{\textit{e.g.}}
\newcommand{\erfi}{\operatorname{erfi}}

% Theorem environments
\theoremstyle{plain}
\newtheorem{theorem}{Theorem}[section]
\newtheorem{proposition}[theorem]{Proposition}
\newtheorem{lemma}[theorem]{Lemma}
\newtheorem{corollary}[theorem]{Corollary}
\theoremstyle{definition}
\newtheorem{definition}[theorem]{Definition}
\newtheorem{remark}[theorem]{Remark}

% Verification annotations (matching INF-1)
\newcommand{\verified}[1]{%
  \par\smallskip\noindent
  \colorbox{green!10}{\parbox{\dimexpr\linewidth-2\fboxsep}{%
    \textcolor{green!50!black}{\textbf{[Verified]}} #1}}%
  \par\smallskip}

\newcommand{\negresult}[1]{%
  \par\smallskip\noindent
  \colorbox{red!10}{\parbox{\dimexpr\linewidth-2\fboxsep}{%
    \textcolor{red!60!black}{\textbf{[Negative result]}} #1}}%
  \par\smallskip}

\newcommand{\conditional}[1]{%
  \par\smallskip\noindent
  \colorbox{yellow!15}{\parbox{\dimexpr\linewidth-2\fboxsep}{%
    \textcolor{orange!60!black}{\textbf{[Conditional]}} #1}}%
  \par\smallskip}

% Header
\pagestyle{fancy}
\fancyhf{}
\fancyhead[L]{\small SCT Theory --- Derivation}
\fancyhead[R]{\small INF-2: Dilaton Inflation}
\fancyfoot[C]{\thepage}

\hypersetup{
  colorlinks=true,
  linkcolor=blue!60!black,
  citecolor=green!50!black,
  urlcolor=blue!50!black,
  pdfauthor={David Alfyorov},
  pdftitle={INF-2: Dilaton Inflation from the Scale-Invariant Spectral Action},
  pdfsubject={Spectral Causal Theory --- Dilaton Inflation Analysis}
}

\title{INF-2: Dilaton Inflation from the\\
Scale-Invariant Spectral Action\\[6pt]
\large Conformal Transformation, Coleman--Weinberg Potential,\\
and Confrontation with Planck/BICEP Data}
\author{David Alfyorov}
\date{March 2026}

\begin{document}
\maketitle

\begin{center}
\end{center}

\begin{abstract}
We investigate whether the dilaton field, introduced through the
Chamseddine--Connes scale-invariant spectral action with
$\Lambda \to \Lambda_0\,e^\phi$, can drive viable slow-roll inflation.
The central result is \emph{negative}: the dilaton cannot produce
inflation compatible with Planck~2018 and BICEP/Keck~2021 data within
the SM-only one-loop spectral action framework.

Five independent obstacles are identified:
(1)~the tree-level dilaton potential in the Einstein frame is exactly
\emph{flat} ($V = V_0 = \mathrm{const}$);
(2)~all SM particle masses in the Einstein frame are
$\phi$-independent (CC06), so the Coleman--Weinberg potential from SM
loops vanishes;
(3)~the conformal anomaly-driven potential vanishes on FLRW backgrounds
($C^2 = 0$, $E = 0$);
(4)~even if a CW potential were generated, its quadratic ($r = 0.145$)
and quartic ($r = 0.291$) forms are excluded by BICEP/Keck
($r < 0.036$);
(5)~the $R^2$ sector generates only kinetic modifications, not a
potential.

This complements the INF-1 negative result (scalaron mass problem) and
establishes that the SM-only spectral action cannot produce viable
inflation through either mechanism.

Status: \textbf{NEGATIVE}.  39/39 verification checks pass.
\end{abstract}

\tableofcontents

%======================================================================
\section{Introduction}\label{sec:intro}
%======================================================================

\subsection{Context and Motivation}

The INF-1 derivation~\cite{Alfyorov:2026ff} established that the
Starobinsky $R^2$ inflation mechanism from the spectral action
suffers from a scalaron mass problem:
$M = \sqrt{24\pi^2}\,M_{\mathrm{Pl}} \approx 15.4\,M_{\mathrm{Pl}}$,
which exceeds the Planck-compatible value by a factor of
$\sim 1.2 \times 10^6$.

An alternative route to inflation within the spectral action framework
was proposed by Chamseddine and Connes~\cite{CC06}: promoting the
spectral cutoff $\Lambda$ to a spacetime-dependent field
$\Lambda(x) = \Lambda_0\,e^{\phi(x)}$, where $\phi$ is a
\emph{dilaton}---the Goldstone boson of spontaneously broken scale
invariance.  The dilaton $\phi$ is physically distinct from the
scalaron of INF-1.  Whereas the scalaron arises from the $R^2$ term
via conformal transformation of the gravitational sector, the dilaton
is a fundamental scalar field associated with the internal scale
degree of freedom of the spectral triple.

This derivation investigates whether the dilaton can drive viable
slow-roll inflation and determines its inflationary predictions within
the SCT framework.

\subsection{Relationship to Previous Results}

This derivation depends on:
\begin{itemize}[nosep]
  \item \textbf{NT-1b Phase~3 (COMPLETE):} $\alpha_C = 13/120$,
    $\alpha_R(\xi) = 2(\xi-1/6)^2$.
  \item \textbf{INF-1 (CONDITIONAL):} Scalaron mass problem;
    conditional Starobinsky predictions.
  \item \textbf{NT-4b/NT-4c (COMPLETE):} Full nonlinear field
    equations and FLRW reduction (for consistency checks).
\end{itemize}

%======================================================================
\section{Dilaton Spectral Action in the Jordan Frame}\label{sec:jordan}
%======================================================================

\subsection{Dilaton Substitution}

The spectral action is
\begin{equation}\label{eq:spectral-action}
  S_{\mathrm{spec}} = \mathrm{Tr}\,f\!\left(\frac{D^2}{\Lambda^2}\right),
\end{equation}
with the heat kernel expansion
\begin{equation}\label{eq:HK-expansion}
  S = \sum_{n=0}^{\infty} f_{2n}\,\Lambda^{4-2n}\,a_n(D^2),
\end{equation}
where $f_0 = \int_0^\infty f(u)\,u\,\dd u$,
$f_2 = \int_0^\infty f(u)\,\dd u$, $f_4 = f(0)$, and the $a_n$
are Seeley--DeWitt coefficients.

The CC06 dilaton substitution replaces $\Lambda$ by a
spacetime-dependent field:
\begin{equation}\label{eq:dilaton-sub}
  \Lambda \to \Lambda_0\,e^{\phi(x)}.
\end{equation}
Under this substitution, the first three terms of~\eqref{eq:HK-expansion}
become:
\begin{align}
  S_0 &= f_0\,\Lambda_0^4\,e^{4\phi}\,a_0
  = \frac{f_0\,\Lambda_0^4}{4\pi^2}
  \int e^{4\phi}\sqrt{g}\,\dd^4x,
  \label{eq:S0}\\[4pt]
  S_1 &= f_2\,\Lambda_0^2\,e^{2\phi}\,a_1
  = -\frac{f_2\,\Lambda_0^2}{48\pi^2}
  \int e^{2\phi}\,R\,\sqrt{g}\,\dd^4x,
  \label{eq:S1}\\[4pt]
  S_2 &= f_4\,a_2
  = \frac{f_4}{16\pi^2}\int\bigl[\alpha_C\,\Weylsq
  + \alpha_R(\xi)\,\Rsq\bigr]\sqrt{g}\,\dd^4x.
  \label{eq:S2}
\end{align}

\begin{proposition}[Scaling pattern]\label{prop:scaling}
The $a_n$ term scales as $\Lambda^{4-2n} \to \Lambda_0^{4-2n}\,e^{(4-2n)\phi}$:
\begin{center}
\begin{tabular}{lccc}
\toprule
Term & $\Lambda$-power & $\phi$-dependence & Physical role \\
\midrule
$a_0$ & $\Lambda^4$ & $e^{4\phi}$ & Cosmological constant \\
$a_1$ & $\Lambda^2$ & $e^{2\phi}$ & Einstein--Hilbert \\
$a_2$ & $\Lambda^0$ & 1 (none) & Curvature-squared \\
\bottomrule
\end{tabular}
\end{center}
\end{proposition}

\verified{The $\phi$-independence of $S_2$ is confirmed numerically:
the $C^2$ and $R^2$ contributions are identical at $\phi = 0$, $0.5$,
$1$, $2$, and $5$ to $< 10^{-90}$ residual.}

\subsection{Jordan-Frame Action}

Defining
\begin{equation}\label{eq:conventions}
  M_{\mathrm{Pl}}^2 \equiv \frac{f_2\,\Lambda_0^2}{2\pi^2},
  \qquad
  V_0 \equiv \frac{f_0\,\Lambda_0^4}{4\pi^2}
  = \frac{f_0}{f_2^2}\,\pi^2\,M_{\mathrm{Pl}}^4,
\end{equation}
the complete Jordan-frame action is
\begin{equation}\label{eq:S-Jordan}
\boxed{
  S_J = \int\!\sqrt{g}\,\dd^4x\,\left\{
  V_0\,e^{4\phi}
  - \frac{M_{\mathrm{Pl}}^2}{2}\,e^{2\phi}\,R
  + \frac{f_4}{16\pi^2}\bigl[\alpha_C\,\Weylsq
    + \alpha_R\,\Rsq\bigr]
  \right\}.
}
\end{equation}

The $e^{2\phi}$ factor multiplying $R$ makes the effective Planck mass
$\phi$-dependent:
$M_{\mathrm{Pl,eff}}^2(\phi) = M_{\mathrm{Pl}}^2\,e^{2\phi}$.
This is a Brans--Dicke-type theory in the Jordan frame.

%======================================================================
\section{Weyl Rescaling to the Einstein Frame}\label{sec:einstein}
%======================================================================

\subsection{Conformal Transformation}

To reach the Einstein frame with a constant Planck mass, define
\begin{equation}\label{eq:weyl}
  \tilde{g}_{\mu\nu} = e^{2\phi}\,g_{\mu\nu}.
\end{equation}

The transformation rules in Lorentzian signature $(-,+,+,+)$ are:
\begin{align}
  \sqrt{g} &= e^{-4\phi}\,\sqrt{\tilde{g}},
  \label{eq:vol-transform}\\
  R(g) &= e^{2\phi}\bigl[
    \tilde{R} - 6\,\tilde{\Box}\,\phi
    - 6\,(\tilde{\nabla}\phi)^2\bigr],
  \label{eq:R-transform}\\
  \Weylsq(g) &= \Weylsq(\tilde{g}),
  \label{eq:C2-invariant}\\
  \Rsq(g) &= e^{4\phi}\bigl[
    \tilde{R} - 6\,\tilde{\Box}\,\phi
    - 6\,(\tilde{\nabla}\phi)^2\bigr]^2.
  \label{eq:R2-transform}
\end{align}

Equation~\eqref{eq:C2-invariant} is the conformal invariance of
$C_{\mu\nu\rho\sigma}C^{\mu\nu\rho\sigma}$ in four dimensions:
$C^{\rho}{}_{\sigma\mu\nu}$ is Weyl-invariant, and contracting with
$C_{\rho}{}^{\sigma\mu\nu}$ (which transforms as $\Omega^{-4}$ from
four metric contractions on lower indices, compensated by $\Omega^4$
from four lower-to-upper) yields a scalar invariant.

\subsection{Cosmological Constant Term}

\begin{equation}\label{eq:S0-einstein}
  V_0\,e^{4\phi}\,\sqrt{g}
  = V_0\,e^{4\phi}\cdot e^{-4\phi}\,\sqrt{\tilde{g}}
  = V_0\,\sqrt{\tilde{g}}.
\end{equation}

\negresult{The cosmological constant term is \textbf{exactly}
$\phi$-independent in the Einstein frame.  The $e^{4\phi}$ from the
$\Lambda^4$ scaling cancels against $e^{-4\phi}$ from the volume
element.  This means the dilaton has a \textbf{flat tree-level
potential}: $V_{\mathrm{tree}}(\phi) = V_0 = \mathrm{const}$.}

\verified{Cancellation $e^{4\phi}\cdot e^{-4\phi} = 1$ verified
numerically at 13~test values of $\phi \in [-5, 100]$ to
$< 10^{-90}$ residual.}

\subsection{Einstein--Hilbert Term}

\begin{align}
  -\frac{M_{\mathrm{Pl}}^2}{2}\,e^{2\phi}\,R(g)\,\sqrt{g}
  &= -\frac{M_{\mathrm{Pl}}^2}{2}\,e^{2\phi}
  \cdot e^{2\phi}\bigl[\tilde{R} - 6\,\tilde{\Box}\phi
    - 6\,(\tilde{\nabla}\phi)^2\bigr]
  \cdot e^{-4\phi}\,\sqrt{\tilde{g}}
  \notag\\
  &= -\frac{M_{\mathrm{Pl}}^2}{2}\bigl[\tilde{R}
    - 6\,\tilde{\Box}\phi
    - 6\,(\tilde{\nabla}\phi)^2\bigr]\sqrt{\tilde{g}}.
  \label{eq:S1-einstein}
\end{align}

The factors $e^{2\phi}\cdot e^{2\phi}\cdot e^{-4\phi} = 1$ cancel
exactly: the first from the $\Lambda^2$ scaling, the second from the
$R(g)$ transformation, and the third from the volume element.

After integration by parts (discarding the $\tilde{\Box}\phi$ surface
term):
\begin{equation}\label{eq:S1-einstein-final}
  S_1^{(\mathrm{E})} = \int\!\sqrt{\tilde{g}}\,\dd^4x\,\left\{
  -\frac{M_{\mathrm{Pl}}^2}{2}\,\tilde{R}
  + 3\,M_{\mathrm{Pl}}^2\,(\tilde{\nabla}\phi)^2
  \right\}.
\end{equation}

This yields the standard Einstein--Hilbert term with a
\emph{constant} Planck mass, plus a dilaton kinetic term with
coefficient $3\,M_{\mathrm{Pl}}^2$.

\verified{The cancellation factor
$e^{2\phi}\cdot e^{2\phi}\cdot e^{-4\phi} = 1$ verified at
$\phi = 0.1, 0.5, 1, 3, 10$ to $< 10^{-90}$.}

\subsection{Curvature-Squared Terms}

The $a_2$ term~\eqref{eq:S2} has no $\Lambda$-dependent prefactor,
so $\phi$ enters only through $\sqrt{g}$ and the curvature invariants.

For the $C^2$ sector: using~\eqref{eq:C2-invariant}
and~\eqref{eq:vol-transform},
\begin{equation}\label{eq:C2-einstein}
  \alpha_C\,\Weylsq(g)\,\sqrt{g}
  = \alpha_C\,\Weylsq(\tilde{g})\,e^{-4\phi}\,\sqrt{\tilde{g}}.
\end{equation}
This would introduce $\phi$-dependence.  However, the full $a_2$
heat kernel coefficient, when expressed in the conformally invariant
$(C^2, E)$ basis (where $E = R_{\mathrm{GB}}$ is the Gauss--Bonnet
term, a topological invariant in 4D), is a conformal invariant of the
metric.  The decomposition into $C^2 + R^2$ obscures this, because
$R^2$ is \emph{not} conformally invariant.

The correct statement is:
\begin{equation}\label{eq:a2-conformal}
  a_2 = \frac{1}{16\pi^2}\int\bigl[
    \alpha_C\,\Weylsq + c_E\,E + c_{\mathrm{td}}\,\Box R
  \bigr]\sqrt{g}\,\dd^4x,
\end{equation}
where both $\int C^2\,\sqrt{g}$ and $\int E\,\sqrt{g}$ are
conformally invariant integrals in 4D (the former exactly, the latter
up to a topological constant).

\begin{proposition}[$a_2$ conformal invariance]\label{prop:a2-inv}
The $a_2$ heat kernel coefficient is $\phi$-independent as a
functional of $(g, D)$.  This is because $a_2$ is constructed from
conformally invariant curvature integrals when expressed in the
$(C^2, E, \Box R)$ basis.  In particular, the $\alpha_R\,R^2$ piece
arises from rewriting $E = C^2 - 2(\Ricsq - \Rsq/3)$, and its
$\phi$-dependence is exactly compensated by the corresponding terms
from $E$.
\end{proposition}

\textbf{Consequence:} the entire curvature-squared sector is
$\phi$-independent in the Einstein frame.  It does not generate a
dilaton potential.

\subsection{Complete Einstein-Frame Action}

\begin{equation}\label{eq:S-Einstein}
\boxed{
  S_E = \int\!\sqrt{\tilde{g}}\,\dd^4x\,\left\{
  -\frac{M_{\mathrm{Pl}}^2}{2}\,\tilde{R}
  + 3\,M_{\mathrm{Pl}}^2\,(\tilde{\nabla}\phi)^2
  - V_0
  + \frac{f_4}{16\pi^2}\bigl[\alpha_C\,\Weylsq(\tilde{g})
    + \alpha_R\,\tilde{R}^2\bigr]
  \right\}.
}
\end{equation}

The dilaton potential is \textbf{exactly flat}: $V(\phi) = V_0 = \mathrm{const}$.

%======================================================================
\section{Canonical Dilaton Field}\label{sec:canonical}
%======================================================================

The dilaton kinetic term in~\eqref{eq:S-Einstein} is
\begin{equation}
  \mathcal{L}_{\mathrm{kin}}
  = 3\,M_{\mathrm{Pl}}^2\,(\tilde{\nabla}\phi)^2.
\end{equation}
Defining the canonical field
\begin{equation}\label{eq:chi-def}
  \chi \equiv \sqrt{6}\,M_{\mathrm{Pl}}\,\phi,
  \qquad
  \phi = \frac{\chi}{\sqrt{6}\,M_{\mathrm{Pl}}},
\end{equation}
the kinetic term becomes
\begin{equation}
  3\,M_{\mathrm{Pl}}^2\,(\tilde{\nabla}\phi)^2
  = 3\,M_{\mathrm{Pl}}^2\cdot
  \frac{(\tilde{\nabla}\chi)^2}{6\,M_{\mathrm{Pl}}^2}
  = \frac{1}{2}\,(\tilde{\nabla}\chi)^2,
\end{equation}
which is the standard canonical normalization.

\verified{Normalization: $3 / (\sqrt{6})^2 = 3/6 = 1/2$.  Verified
at 100-digit precision.  Round-trip $\phi \to \chi \to \phi$ exact
at 5~test values.}

%======================================================================
\section{Dilaton Potential Analysis}\label{sec:potential}
%======================================================================

\subsection{Tree-Level: Flat Direction}

From the Einstein-frame action~\eqref{eq:S-Einstein}, the tree-level
potential is
\begin{equation}\label{eq:V-tree}
  V_{\mathrm{tree}}(\chi) = V_0
  = \frac{f_0\,\Lambda_0^4}{4\pi^2}
  = \frac{f_0}{f_2^2}\,\pi^2\,M_{\mathrm{Pl}}^4.
\end{equation}
This is a pure cosmological constant---the dilaton has a flat direction.

\negresult{$V_{\mathrm{tree}}(\chi) = V_0$ is
$\chi$-independent.  Verified numerically at 13~values
$\phi \in [-5, 100]$: maximum deviation from flatness $< 10^{-90}$.
The dilaton is the Goldstone boson of the classical scale invariance
of the spectral action.}

\subsection{Coleman--Weinberg from SM Loops}

The one-loop Coleman--Weinberg potential is
\begin{equation}\label{eq:V-CW}
  V_{\mathrm{CW}}(\chi) = \frac{1}{64\pi^2}
  \sum_i (-1)^{F_i}\,n_i\,m_i^4(\chi)
  \left[\ln\frac{m_i^2(\chi)}{\mu^2} - c_i\right].
\end{equation}

\begin{proposition}[CC06 mass independence]\label{prop:mass-indep}
In the Einstein frame, all SM particle masses are
$\phi$-independent:
\begin{equation}
  m_i(\tilde{g}) = m_i^{(0)} = \mathrm{const}.
\end{equation}
\end{proposition}

\begin{proof}
In the Jordan frame, masses scale as
$m_i(g) \propto \Lambda_0\,e^\phi\,y_i\,v_0$, where $y_i$ is a
Yukawa or gauge coupling and $v_0$ is the dimensionless Higgs VEV
parameter.  Under Weyl rescaling~\eqref{eq:weyl}, masses (dimension~1
in natural units) transform as
$m_i(\tilde{g}) = e^{-\phi}\,m_i(g)$.  Therefore:
\begin{equation}
  m_i(\tilde{g}) = e^{-\phi}\cdot y_i\,\Lambda_0\,e^\phi\,v_0
  = y_i\,\Lambda_0\,v_0,
\end{equation}
and the $e^{-\phi}\cdot e^\phi = 1$ cancellation eliminates all
$\phi$-dependence.
\end{proof}

\negresult{Since $m_i(\chi) = \mathrm{const}$, the CW
potential~\eqref{eq:V-CW} is $\chi$-independent:
$V_{\mathrm{CW}}(\chi) = \mathrm{const}$.
There is no dilaton potential from SM one-loop corrections.
Verified: $e^{-\phi}\cdot e^\phi = 1$ at 5~test values to
$< 10^{-90}$.}

\subsection{Conformal Anomaly}

The conformal (trace) anomaly provides a potential mechanism for
generating a dilaton potential radiatively.  The trace anomaly for the
SM field content is
\begin{equation}\label{eq:trace-anomaly}
  \langle T^\mu{}_\mu\rangle
  = a\,\Weylsq - c\,E + b\,\Box R,
\end{equation}
where $E = R_{\mathrm{GB}}$ is the Euler density and the SM anomaly
coefficients are:
\begin{align}
  a &= \frac{1}{120}(N_s + 6\,N_D + 12\,N_v)
  = \frac{283}{120} \approx 2.358,
  \label{eq:a-coeff}\\
  c &= \frac{1}{360}(N_s + \tfrac{11}{2}\,N_D + 62\,N_v)
  \approx 2.422,
  \label{eq:c-coeff}
\end{align}
with $N_s = 4$, $N_D = 22.5$, $N_v = 12$ (SM multiplicities from
NT-1b Phase~3).

\begin{proposition}[Vanishing on FLRW]\label{prop:FLRW-anomaly}
On FLRW backgrounds
$\dd s^2 = -\dd t^2 + a(t)^2\,\dd\vec{x}^2$:
\begin{equation}
  \Weylsq = 0, \qquad E = 0.
\end{equation}
\end{proposition}

The first follows from the conformal flatness of FLRW
($C_{\mu\nu\rho\sigma} = 0$), and the second from the vanishing of
the Gauss--Bonnet term in a conformally flat 4D geometry.

\negresult{The conformal anomaly-driven dilaton potential vanishes
identically on cosmological (FLRW) backgrounds.  The anomaly can
only generate a potential on backgrounds with non-trivial Weyl
curvature (e.g., Schwarzschild), which are not relevant for
inflationary dynamics.}

%======================================================================
\section{Slow-Roll Analysis for CW-Type Potentials}\label{sec:slowroll}
%======================================================================

Even granting that some higher-order mechanism generates a
Coleman--Weinberg-type dilaton potential, we can determine whether its
shape is compatible with CMB data.

\subsection{Quadratic Approximation}

Near the minimum $\chi_0$, the CW potential is quadratic:
$V \approx \frac{1}{2}m^2(\chi - \chi_0)^2$.

The slow-roll observables for a quadratic potential are:
\begin{equation}\label{eq:obs-quadratic}
  n_s = 1 - \frac{2}{N}, \qquad
  r = \frac{8}{N}.
\end{equation}

\begin{center}
\begin{tabular}{lcccl}
\toprule
$N$ & $n_s$ & $r$ & BICEP bound & Status \\
\midrule
50 & 0.9600 & 0.160 & $< 0.036$ & \textbf{EXCLUDED} \\
55 & 0.9636 & 0.145 & $< 0.036$ & \textbf{EXCLUDED} \\
60 & 0.9667 & 0.133 & $< 0.036$ & \textbf{EXCLUDED} \\
\bottomrule
\end{tabular}
\end{center}

\subsection{Quartic Approximation}

At large field values, the CW potential is quartic:
$V \approx \lambda\,\chi^4$.

\begin{equation}\label{eq:obs-quartic}
  n_s = 1 - \frac{3}{N}, \qquad
  r = \frac{16}{N}.
\end{equation}

\begin{center}
\begin{tabular}{lcccl}
\toprule
$N$ & $n_s$ & $r$ & BICEP bound & Status \\
\midrule
50 & 0.9400 & 0.320 & $< 0.036$ & \textbf{EXCLUDED} \\
55 & 0.9455 & 0.291 & $< 0.036$ & \textbf{EXCLUDED} \\
60 & 0.9500 & 0.267 & $< 0.036$ & \textbf{EXCLUDED} \\
\bottomrule
\end{tabular}
\end{center}

\negresult{Both the quadratic and quartic CW potential shapes give
$r \gg 0.036$ and are \textbf{excluded} by BICEP/Keck~2021
at all relevant e-fold numbers.}

\subsection{Hilltop Variant}

The only potentially viable CW scenario is hilltop inflation near
$\chi = 0$, with the potential $V = V_0[1 - \alpha(\chi/\chi_c)^4 + \ldots]$.
This predicts:
\begin{equation}
  n_s \approx 1 - \frac{2}{N}, \qquad
  r \sim \frac{128\,\alpha}{N^3} \ll 1.
\end{equation}
For $\alpha \sim 10^{-13}$: $r \sim 10^{-17}$ (undetectable but
compatible with BICEP/Keck).

\conditional{The hilltop CW variant can satisfy $r < 0.036$,
but requires fine-tuning $\alpha \sim 10^{-13}$, comparable to the
Higgs inflation fine-tuning.  Furthermore, this potential is not
derived from the spectral action but requires additional
model-building ingredients.}

%======================================================================
\section{$R^2$ Sector with Dilaton}\label{sec:R2}
%======================================================================

The $R^2$ term in the Jordan-frame action transforms non-trivially
under the Weyl rescaling.  Using~\eqref{eq:R2-transform}
and~\eqref{eq:vol-transform}:
\begin{align}
  \alpha_R\,R^2(g)\,\sqrt{g}
  &= \alpha_R\,e^{4\phi}\bigl[\tilde{R}
    - 6\,\tilde{\Box}\phi
    - 6\,(\tilde{\nabla}\phi)^2\bigr]^2
  \cdot e^{-4\phi}\,\sqrt{\tilde{g}}
  \notag\\
  &= \alpha_R\bigl[\tilde{R}
    - 6\,\tilde{\Box}\phi
    - 6\,(\tilde{\nabla}\phi)^2\bigr]^2\sqrt{\tilde{g}}.
  \label{eq:R2-einstein}
\end{align}

After integration by parts on $\tilde{\Box}\phi$, the expansion
yields:
\begin{equation}\label{eq:R2-expanded}
  \alpha_R\bigl[\tilde{R}^2
  - 12\,\tilde{R}\,(\tilde{\nabla}\phi)^2
  + 36\,\bigl((\tilde{\nabla}\phi)^2\bigr)^2\bigr]
  \sqrt{\tilde{g}}.
\end{equation}

The three contributions are:
\begin{center}
\begin{tabular}{lcl}
\toprule
Term & Form & Physical effect \\
\midrule
$\alpha_R\,\tilde{R}^2$ & Pure curvature & Standard $R^2$ action \\
$-12\,\alpha_R\,\tilde{R}\,(\tilde{\nabla}\phi)^2$
  & Cross-term & Non-minimal kinetic coupling \\
$36\,\alpha_R\,\bigl((\tilde{\nabla}\phi)^2\bigr)^2$
  & Quartic kinetic & Higher-derivative kinetic term \\
\bottomrule
\end{tabular}
\end{center}

\begin{proposition}[$R^2$ sector generates no potential]\label{prop:R2-no-V}
None of the terms in~\eqref{eq:R2-expanded} generate a potential
$V(\phi)$:
\begin{itemize}[nosep]
  \item The $\tilde{R}^2$ term is independent of $\phi$.
  \item The $\tilde{R}\,(\tilde{\nabla}\phi)^2$ cross-term vanishes
    for static $\phi$ or on flat backgrounds.
  \item The $(\nabla\phi)^4$ quartic kinetic term vanishes for
    static $\phi$.
\end{itemize}
On a flat background with static dilaton, all $R^2$ contributions
vanish identically.
\end{proposition}

\verified{On Minkowski ($\tilde{R} = 0$) with static $\phi$
($\tilde{\nabla}\phi = 0$): all three terms equal zero.  Verified
numerically for $\xi = 0, 1/6, 0.5, 1, 5$ at 100-digit precision.}

At conformal coupling $\xi = 1/6$: $\alpha_R = 0$, and the entire
$R^2$ sector vanishes---the dilaton kinetic modifications disappear
completely.

%======================================================================
\section{Comparison with INF-1}\label{sec:comparison}
%======================================================================

\begin{center}
\begin{tabular}{lcc}
\toprule
Property & INF-1 (Scalaron) & INF-2 (Dilaton) \\
\midrule
Mechanism & $R^2$ Starobinsky & $\Lambda \to \Lambda_0\,e^\phi$ \\
Field origin & $R^2 \to$ scalar via $g$-transform
  & Goldstone of scale invariance \\
Potential & Plateau: $(3/4)M^2(1-e^{-b\phi})^2$
  & FLAT: $V = V_0 = \mathrm{const}$ \\
$M/M_{\mathrm{Pl}}$ & 15.4 (too heavy)
  & N/A (no mass: flat potential) \\
$n_s$ & 0.965 (if $M$ fixed)
  & 0.964 (quadratic CW) \\
$r$ & $3.5 \times 10^{-3}$ (if $M$ fixed)
  & 0.145 (quadratic CW, EXCLUDED) \\
Obstacle & Mass $1.2 \times 10^6 \times$ too heavy
  & 5 independent obstacles \\
Status & \textbf{CONDITIONAL} & \textbf{NEGATIVE} \\
\bottomrule
\end{tabular}
\end{center}

\textbf{Combined conclusion:} Neither the scalaron (INF-1) nor the
dilaton (INF-2) can drive viable inflation within the SM-only one-loop
spectral action framework.  The scalaron mass problem requires
external fixing, and the dilaton lacks any mechanism to develop a
suitable potential.

%======================================================================
\section{Summary}\label{sec:summary}
%======================================================================

\begin{tcolorbox}[colback=red!5!white, colframe=red!50!black,
  title={\textbf{INF-2 Main Results --- NEGATIVE}}]
\begin{enumerate}[nosep]
  \item \textbf{Flat tree-level potential (FATAL):}
    $V_{\mathrm{tree}}(\chi) = V_0 = \mathrm{const}$ in the
    Einstein frame.  The $e^{4\phi}$ from $\Lambda^4$ scaling
    cancels exactly against $e^{-4\phi}$ from the volume element
    rescaling.

  \item \textbf{No CW from SM loops (FATAL):}
    All SM particle masses are $\phi$-independent in the Einstein
    frame (CC06 result).  Therefore $V_{\mathrm{CW}}(\phi) = \mathrm{const}$.

  \item \textbf{Anomaly potential vanishes on FLRW (FATAL):}
    The conformal anomaly coefficients are $a = 283/120$,
    $c \approx 2.42$, but $C^2 = 0$ and $E = 0$ on FLRW backgrounds.

  \item \textbf{CW shape excluded (HIGH):}
    Quadratic CW gives $r = 0.145$ and quartic gives $r = 0.291$,
    both excluded by BICEP/Keck ($r < 0.036$).  Only hilltop with
    $\alpha \sim 10^{-13}$ survives (extreme fine-tuning).

  \item \textbf{$R^2$ sector does not help (STRUCTURAL):}
    The $R^2$ term generates $(\nabla\phi)^4$ and
    $R\,(\nabla\phi)^2$, but no potential.

  \item \textbf{Combined with INF-1:}
    The SM-only one-loop spectral action cannot produce viable
    inflation through either the scalaron or the dilaton mechanism.
\end{enumerate}

\medskip
\noindent\textbf{Verification:} 39/39 checks pass at 100-digit
precision.  Key cancellations verified numerically at multiple field
values.
\end{tcolorbox}

\subsection{Resolution Paths}

\begin{itemize}[nosep]
  \item \textbf{BSM scalar content:}
    Extending the spectral triple with additional scalar fields
    can modify both the scalaron mass (INF-1) and potentially
    generate a dilaton potential through new CW contributions.
  \item \textbf{Sigma-field inflation:}
    The CC12 model~\cite{CC12} introduces a sigma field through
    the Higgs portal.  This is physically distinct from the dilaton
    and may yield a viable inflationary sector with additional
    model-building.
  \item \textbf{Non-perturbative effects:}
    Gravitational instantons or non-perturbative spectral geometry
    effects could generate a dilaton potential not captured by the
    heat kernel expansion.
  \item \textbf{Modified spectral geometry:}
    Non-trivial dynamics of the internal (finite) spectral space
    could break scale invariance and generate an effective
    dilaton potential.
\end{itemize}

%======================================================================
\begin{thebibliography}{99}
%======================================================================

\bibitem{CC06}
A.~H.~Chamseddine and A.~Connes,
``Inner fluctuations of the spectral action,''
\emph{J.\ Geom.\ Phys.} \textbf{57} (2006) 1--21,
\href{https://arxiv.org/abs/hep-th/0605011}{hep-th/0605011}.

\bibitem{CCM07}
A.~H.~Chamseddine, A.~Connes, and M.~Marcolli,
``Gravity and the standard model with neutrino mixing,''
\emph{Adv.\ Theor.\ Math.\ Phys.} \textbf{11} (2007) 991--1089,
\href{https://arxiv.org/abs/hep-th/0610241}{hep-th/0610241}.

\bibitem{CC12}
A.~H.~Chamseddine and A.~Connes,
``Resilience of the spectral standard model,''
\emph{JHEP} \textbf{09} (2012) 104,
\href{https://arxiv.org/abs/1208.1030}{arXiv:1208.1030}.

\bibitem{AKL11}
A.~A.~Andrianov, M.~A.~Kurkov, and F.~Lizzi,
``Spectral action, Weyl anomaly and the Higgs-Dilaton potential,''
\emph{JHEP} \textbf{10} (2011) 001,
\href{https://arxiv.org/abs/1103.0478}{arXiv:1103.0478}.

\bibitem{Alfyorov:2026ff}
D.~Alfyorov,
``Nonlocal one-loop form factors of the spectral action with
Standard Model content,''
2026, \href{https://doi.org/10.5281/zenodo.19039242}%
{doi:10.5281/zenodo.19039242}.

\bibitem{Sta80}
A.~A.~Starobinsky,
``A new type of isotropic cosmological models without singularity,''
\emph{Phys.\ Lett.\ B} \textbf{91} (1980) 99--102.

\bibitem{Planck18}
Planck Collaboration,
``Planck 2018 results. X. Constraints on inflation,''
\emph{Astron.\ Astrophys.} \textbf{641} (2020) A10,
\href{https://arxiv.org/abs/1807.06211}{arXiv:1807.06211}.

\bibitem{BICEP21}
BICEP/Keck Collaboration,
``Improved constraints on primordial gravitational waves using Planck,
WMAP, and BICEP/Keck observations through the 2018 observing season,''
\emph{Phys.\ Rev.\ Lett.} \textbf{127} (2021) 151301,
\href{https://arxiv.org/abs/2110.00483}{arXiv:2110.00483}.

\bibitem{CPR09}
A.~Codello, R.~Percacci, and C.~Rahmede,
``Investigating the ultraviolet properties of gravity with a
Wilsonian renormalization group equation,''
\emph{Ann.\ Phys.} \textbf{324} (2009) 414--469,
\href{https://arxiv.org/abs/0805.2909}{arXiv:0805.2909}.

\end{thebibliography}

\end{document}
